\documentclass[11pt]{article}
\usepackage[T1]{fontenc}
\usepackage{amsmath,amssymb,amsthm}
\usepackage[margin=1in]{geometry}
\usepackage[hidelinks]{hyperref}
\hypersetup{pdftitle={All Deformation Maps in Example 8.7 of Extending Structures for n-Lie Algebras},pdfauthor={Xiaojian Zeng}}

\newtheorem{theorem}{Theorem}
\newcommand{\K}{\mathbb K}

\title{All Deformation Maps in Example 8.7 of\\
\emph{Extending Structures for $n$-Lie Algebras}}
\author{Xiaojian Zeng\\
PBC School of Finance, Tsinghua University, 100190\\
\texttt{xiaojianzeng@foxmail.com}}
\date{}

\begin{document}
\maketitle

\begin{abstract}
Let $n\geq4$ and let $\K$ be a field of characteristic zero.  For the
$n$-Lie matched pair in Example 8.7 of Zhang's \emph{Extending Structures
for $n$-Lie Algebras}, we classify every linear deformation map
$r\colon V\to\mathfrak g$.  In the column convention
\[
R=\begin{pmatrix}A&v\\ w^{\mathsf T}&\alpha\end{pmatrix},
\]
the complete criterion is
$\det(A)v=0$ and $\det(R)=\alpha\det(A)$.  Thus, if $A$ is invertible,
then $v=0$ and $w,\alpha$ are arbitrary; if $A$ is singular, the sole
condition is $\det(R)=0$.  The proof reduces the deformation equation on
arbitrary inputs to two determinant-valued alternating maps and includes an
explicit audit of the Laplace-expansion sign.
\end{abstract}

\section{Context and statement}

Filippov introduced $n$-Lie algebras as algebras with an alternating
$n$-linear bracket satisfying the fundamental identity
\cite{Filippov1985}.  Zhang developed extending structures and deformation
maps first for $3$-Lie algebras \cite{Zhang2022} and subsequently for
$n$-Lie algebras \cite{Zhang2026}.  The latter paper's Example 8.7 gives a
matched pair for every $n>3$, displays one deformation map, and uses it in
the subsequent classification of complement isomorphism types.  Here we
answer the separate all-matrices question for that example.  We make no
novelty or priority claim relative to specializations not compared in full.

Fix an integer $n\geq4$ and a characteristic-zero field $\K$.  Let
\[
 \mathfrak g=\operatorname{span}_{\K}\{x_1,\ldots,x_n\},\qquad
 V=\operatorname{span}_{\K}\{u_1,\ldots,u_n\}.
\]
The only nonzero ordered basis value of the bracket of $\mathfrak g$ is
\[
 [x_1,\ldots,x_n]_{\mathfrak g}=x_n,
\]
and $V$ is abelian.  The only nonzero ordered basis value of the action
$\rho_{\mathfrak g}\colon\Lambda^{n-1}\mathfrak g\to
\operatorname{End}(V)$ is
\[
 \rho_{\mathfrak g}(x_1,\ldots,x_{n-1})u_n=u_n.
\]
All other matched-pair structure maps in the example vanish.  Consequently,
the deformation-map condition in Definition 7.1 of \cite{Zhang2026}
specializes to
\begin{equation}\label{eq:DM}
 [r(q_1),\ldots,r(q_n)]_{\mathfrak g}
 =r\!\left(\sum_{i=1}^n(-1)^{n-i}
 \rho_{\mathfrak g}\bigl(r(q_1),\ldots,\widehat{r(q_i)},\ldots,
 r(q_n)\bigr)q_i\right)
\end{equation}
for every $q_1,\ldots,q_n\in V$.

Use the column convention
\[
 r(u_j)=\sum_{i=1}^n R_{ij}x_i,
 \qquad
 R=\begin{pmatrix}A&v\\ w^{\mathsf T}&\alpha\end{pmatrix},
\]
where $A\in M_{n-1}(\K)$, $v,w\in\K^{n-1}$ and $\alpha\in\K$.

\begin{theorem}\label{thm:classification}
The map $r$ satisfies \eqref{eq:DM} if and only if
\begin{equation}\label{eq:criterion}
 \det(A)v=0,
 \qquad
 \det(R)=\alpha\det(A).
\end{equation}
Equivalently, the complete solution set has the following two disjoint
parts:
\begin{enumerate}
 \item if $\det(A)\ne0$, then $v=0$, while $w$ and $\alpha$ are arbitrary;
 \item if $\det(A)=0$, then the sole condition is $\det(R)=0$.
\end{enumerate}
\end{theorem}

\section{Determinant reduction and proof}

For $1\leq a\leq n$, let
$p_a\colon\mathfrak g\to\K$ be the $x_a$-coordinate functional, and put
$\varepsilon=u_n^*\colon V\to\K$.  Alternation and multilinearity of the
displayed basis values give
\begin{align}
 [y_1,\ldots,y_n]_{\mathfrak g}
 &=\det\bigl(p_a(y_b)\bigr)_{1\leq a,b\leq n}\,x_n,
 \label{eq:bracketdet}\\
 \rho_{\mathfrak g}(y_1,\ldots,y_{n-1})q
 &=\det\bigl(p_a(y_b)\bigr)_{1\leq a,b\leq n-1}
   \varepsilon(q)u_n.
 \label{eq:actiondet}
\end{align}

Let $q_1,\ldots,q_n\in V$, and let $Q$ be the matrix whose $j$th column
is the coordinate vector of $q_j$ in the ordered $u$-basis.  The coordinate
matrix of $r(q_1),\ldots,r(q_n)$ is $RQ$.  Hence \eqref{eq:bracketdet}
gives
\begin{equation}\label{eq:G}
 [r(q_1),\ldots,r(q_n)]_{\mathfrak g}
 =\det(R)\det(Q)x_n.
\end{equation}

It remains to compute the vector inside $r$ on the right of
\eqref{eq:DM}.
Form
\[
 B=
 \begin{pmatrix}
 p_1(r(q_1))&\cdots&p_1(r(q_n))\\
 \vdots&&\vdots\\
 p_{n-1}(r(q_1))&\cdots&p_{n-1}(r(q_n))\\
 \varepsilon(q_1)&\cdots&\varepsilon(q_n)
 \end{pmatrix}.
\]
After substituting \eqref{eq:actiondet}, the scalar multiplying $u_n$ in
that vector is
\[
 \sum_{i=1}^n(-1)^{n-i}\varepsilon(q_i)
 \det(B_{\widehat n,\widehat i}).
\]
This is precisely the last-row Laplace expansion of $\det B$: its cofactor
sign is $(-1)^{n+i}$, which equals $(-1)^{n-i}$ because the exponents differ
by $2i$.  Moreover,
\[
 B=DQ,
 \qquad
 D=\begin{pmatrix}A&v\\0&1\end{pmatrix},
 \qquad
 \det D=\det A.
\]
Therefore the complete signed action sum is
\begin{equation}\label{eq:H}
 \sum_{i=1}^n(-1)^{n-i}
 \rho_{\mathfrak g}\bigl(r(q_1),\ldots,\widehat{r(q_i)},\ldots,
 r(q_n)\bigr)q_i
 =\det(A)\det(Q)u_n.
\end{equation}

\begin{proof}[Proof of Theorem~\ref{thm:classification}]
Equations \eqref{eq:G} and \eqref{eq:H} show that the residual in
\eqref{eq:DM} on an arbitrary tuple is
\[
 \det(Q)\bigl(\det(R)x_n-\det(A)r(u_n)\bigr).
\]
Thus \eqref{eq:DM} holds for every tuple if and only if it holds for the
ordered basis tuple, for which $Q=I_n$.  Since
\[
 r(u_n)=\sum_{j=1}^{n-1}v_jx_j+\alpha x_n,
\]
the resulting vector equation
$\det(R)x_n=\det(A)r(u_n)$ is equivalent, coordinate by coordinate, to
\eqref{eq:criterion}.  This proves necessity and sufficiency without
cancelling either determinant.

If $\det A\ne0$, the first equation forces $v=0$.  The matrix $R$ is then
block lower triangular, so the second equation holds automatically for
arbitrary $w$ and $\alpha$.  If $\det A=0$, the first equation is automatic
and the second is exactly $\det R=0$.  These cases are exhaustive.
\end{proof}

\section{Exact finite check and scope}

For $n=4$, a supplementary SymPy \cite{MeurerEtAl2017} calculation expands
the specialized equation over $\mathbb Z[r_{ij}]$.  With $A$ the upper-left
$3\times3$ block, its four residual coordinates are exactly
\[
 -\det(A)r_{14},\quad -\det(A)r_{24},\quad -\det(A)r_{34},\quad
 \det(R)-r_{44}\det(A).
\]
This is an exact sign and coordinate check, not a substitute for the
arbitrary-$n$ proof.  Exact rational matrix certificates provide a related
example of keeping finite algebraic verification distinct from broader
claims \cite{ZengLiuRatnavelu2026}; the present verification instead checks
a polynomial identity and the theorem follows from the determinant argument
above.

The theorem retains the source's scope $n\geq4$ over characteristic-zero
fields.  Zhang's earlier paper establishes the ternary extending-structure
and deformation-map framework \cite{Zhang2022}, but no claim is made here
about whether a ternary matrix specialization equivalent to
Theorem~\ref{thm:classification} has appeared elsewhere.

\section*{Computational and AI assistance disclosure}

OpenAI Codex assisted with source inspection, proof derivation, adversarial
sign checking, exact symbolic verification, and manuscript preparation.
The finite check was run in Python with SymPy and used exact symbolic
arithmetic; no floating-point evidence is used in the proof.  The submitting
author remains responsible for verifying every statement and for complying
with the disclosure and authorship policies of any venue.

\bibliographystyle{abbrv}
\bibliography{references}

\end{document}
